

\documentclass[12pt,a4paper,oneside]{book}

% --------------------------------------------------
% Codificação, língua e fontes
% --------------------------------------------------
\usepackage[utf8]{inputenc}
\usepackage[T1]{fontenc}
\usepackage{lmodern}
\usepackage[portuguese]{babel}

% --------------------------------------------------
% Layout e microtipografia
% --------------------------------------------------
\usepackage[a4paper,margin=2.5cm]{geometry}
\usepackage{microtype}
\usepackage{setspace}

% --------------------------------------------------
% Matemática e ambientes de teorema
% --------------------------------------------------
\usepackage{amsmath,amssymb,amsthm}

% --------------------------------------------------
% Gráficos, cores
% --------------------------------------------------
\usepackage{graphicx}
\usepackage{xcolor}

% --------------------------------------------------
% Hiperligações e referências cruzadas
% --------------------------------------------------
\usepackage{hyperref}
\hypersetup{
colorlinks=true,
linkcolor=blue,
citecolor=blue,
urlcolor=blue,
pdftitle={O Código Oculto da Inteligência Artificial --- Terceira edição em desenvolvimento},
pdfauthor={Mário Prazeres}
}
\usepackage[nameinlink,noabbrev]{cleveref}

% --------------------------------------------------
% Listas, código, citações
% --------------------------------------------------
\usepackage{enumitem}
\usepackage{listings}
\usepackage{csquotes}
\usepackage[numbers,sort&compress]{natbib}
\usepackage{booktabs,tabularx,longtable,array}
\usepackage{tikz}
\usetikzlibrary{arrows.meta,positioning,calc}
\usepackage{xurl}
\setcounter{tocdepth}{1}
\setlength{\emergencystretch}{2em}
\InputIfFileExists{revisao.tex}{}{}
\providecommand{\revisao}{trabalho-v3-2026-11-04}
\newcommand{\ficcao}{\textbf{Caso fictício.} }


% Configuração básica para listagens de código
\lstset{
basicstyle=\ttfamily\footnotesize,
  upquote=true,
  breakatwhitespace=false,
  xleftmargin=0pt,
  xrightmargin=0pt,
breaklines=true,
columns=fullflexible,
frame=single,
keepspaces=true,
showstringspaces=false,
tabsize=2,
literate={á}{{\'a}}1 {à}{{\`a}}1 {â}{{\^a}}1 {ã}{{\~a}}1
{é}{{\'e}}1 {ê}{{\^e}}1 {í}{{\'i}}1 {ó}{{\'o}}1 {ô}{{\^o}}1
{õ}{{\~o}}1 {ú}{{\'u}}1 {ü}{{\"u}}1 {ç}{{\c{c}}}1
{Á}{{\'A}}1 {À}{{\`A}}1 {Â}{{\^A}}1 {Ã}{{\~A}}1
{É}{{\'E}}1 {Ê}{{\^E}}1 {Í}{{\'I}}1 {Ó}{{\'O}}1 {Ô}{{\^O}}1
{Õ}{{\~O}}1 {Ú}{{\'U}}1 {Ç}{{\c{C}}}1
{“}{{``}}1 {”}{{''}}1 {‘}{{`}}1 {’}{{'}}1
{–}{{--}}1 {—}{{---}}1 {…}{{...}}3 {→}{{$\to$}}2 {⇒}{{$\Rightarrow$}}2
}

% --------------------------------------------------
% Caixas especiais (Prompt Labs, avisos, etc.)
% --------------------------------------------------
\usepackage[most]{tcolorbox}
\tcbset{
enhanced,
sharp corners,
boxrule=0.6pt,
colback=white,
colframe=black
}

% --------------------------------------------------
% Ambientes de teorema (para quando forem úteis)
% --------------------------------------------------
\theoremstyle{definition}
\newtheorem{definition}{Definição}[chapter]
\theoremstyle{plain}
\newtheorem{theorem}{Teorema}[chapter]
\newtheorem{proposition}{Proposição}[chapter]
\theoremstyle{remark}
\newtheorem{remark}{Nota}[chapter]
\newtheorem{example}{Exemplo}[chapter]

% --------------------------------------------------
% Contadores para elementos especiais
% --------------------------------------------------
\newcounter{promptlab}
\newcounter{godmode}
\newcounter{checklist}
\newcounter{warningbox}
\newcounter{keyidea}

% --------------------------------------------------
% Ambiente: Prompt Lab
% --------------------------------------------------
\newtcolorbox{promptlabbox}[2][]{
breakable,
title={Prompt Lab~\refstepcounter{promptlab}\thepromptlab: #2},
fonttitle=\bfseries,
#1
}

% --------------------------------------------------
% Ambiente: God Mode
% --------------------------------------------------
\newtcolorbox{godmodebox}[2][]{
breakable,
title={God Mode~\refstepcounter{godmode}\thegodmode: #2},
fonttitle=\bfseries,
#1
}

% --------------------------------------------------
% Ambiente: Checklist
% --------------------------------------------------
\newtcolorbox{checklistbox}[2][]{
breakable,
title={Checklist~\refstepcounter{checklist}\thechecklist: #2},
fonttitle=\bfseries,
#1
}

% --------------------------------------------------
% Ambiente: Aten\c c~ao / Aviso
% --------------------------------------------------
\newtcolorbox{warningboxenv}[2][]{
breakable,
title={Atenção \refstepcounter{warningbox}\thewarningbox: #2},
fonttitle=\bfseries,
#1
}

% --------------------------------------------------
% Ambiente: Ideia-Chave
% --------------------------------------------------
\newtcolorbox{keyideabox}[2][]{
breakable,
title={Ideia-chave~\refstepcounter{keyidea}\thekeyidea: #2},
fonttitle=\bfseries,
#1
}

% --------------------------------------------------
% Título, autor, data
% --------------------------------------------------
